% Identifiability and the analytic-vs-rollout gap in bilevel IOC.
% Intended as a citable/includable note for the IOC paper writeup.
% Corresponds to ioc/analytic.py (kkt_fit, cioc_fit) and ioc/inner.py (implicit adjoint).

\documentclass[11pt]{article}
\usepackage{amsmath,amssymb,amsthm}
\usepackage[margin=1in]{geometry}
\usepackage{hyperref}

\newtheorem{proposition}{Proposition}
\newtheorem{remark}{Remark}

\newcommand{\R}{\mathbb{R}}
\newcommand{\theovec}{\theta}
\newcommand{\xdemo}{\tilde{x}}

\title{Identifiability, and why analytic and rollout-based IOC baselines\\
fail differently}
\author{}
\date{}

\begin{document}
\maketitle

\section{Setup}

Let a scene context $c \in \mathcal{C}$ index a trajectory optimization problem
with decision variable $x \in \R^{n_x}$ and a cost
\begin{equation}
  J_\theovec(x, c) \;=\; \sum_{k=1}^{K} \theovec_k\, \phi_k(x, c),
  \qquad \theovec \in \Delta^{K-1},
\end{equation}
linear in a feature basis $\phi = (\phi_1,\dots,\phi_K)$ and a weight vector
$\theovec$ constrained to the simplex. The \emph{inner problem} for context $c$
is
\begin{equation}
  x^\star(\theovec, c) \;=\; \arg\min_x J_\theovec(x, c)
  \quad\text{s.t.}\quad
  F\big(x^\star(\theovec,c),\, \theovec,\, c\big) = 0,
\end{equation}
where $F = \nabla_x J_\theovec$ is the first-order stationarity (KKT) residual
of the inner problem (equality/inequality constraints on $x$, e.g.\ dynamics
and collision, are folded into $F$ and suppressed here for notation).

A demonstration set $\mathcal{D} = \{(c_i, \xdemo_i)\}_{i=1}^{M}$ is assumed
generated by some unknown $\theovec^\star$, with $\xdemo_i$ either an exact
solve $x^\star(\theovec^\star, c_i)$ or a noisy perturbation of one. The IOC
problem is to recover $\theovec^\star$ from $\mathcal{D}$.

\section{Two families of estimator}

\subsection{Analytic (residual-only) estimators}

Define the per-context feature-gradient matrix
\begin{equation}
  B(c) \;=\; \big[\nabla_x \phi_1(\xdemo, c) \;\cdots\; \nabla_x \phi_K(\xdemo,
  c)\big] \in \R^{n_x \times K},
\end{equation}
evaluated \emph{at the demonstration itself}, so that
$F(\xdemo, \theovec, c) = B(c)\,\theovec$.
Inverse KKT (Keshavarz et al., 2011; Englert et al., 2017) fits
\begin{equation}
  \hat\theovec_{\mathrm{KKT}}
  \;=\; \arg\min_{\theovec \in \Delta^{K-1}} \;
  \frac{1}{M}\sum_{i=1}^M \big\| B(c_i)\,\theovec \big\|_2^2
  \;=\; \arg\min_{\theovec \in \Delta^{K-1}} \theovec^\top G\, \theovec,
  \qquad
  G \;=\; \frac{1}{M}\sum_{i=1}^M B(c_i)^\top B(c_i).
  \label{eq:kkt}
\end{equation}
No inner solve of any kind appears in \eqref{eq:kkt}: the estimator only
requires evaluating $\nabla_x \phi_k$ at the given points $\xdemo_i$, which is
why it costs zero forward trajectory-optimization solves.

Continuous IOC (Levine \& Koltun, 2012) is the same construction one order
higher: under a Boltzmann model $p(x) \propto \exp(-J_\theovec(x,c))$, a
second-order (Laplace) expansion of $J_\theovec$ about $\xdemo$ gives, per
context,
\begin{equation}
  -\log p(\xdemo \mid \theovec, c) \;\approx\;
  \tfrac{1}{2}\, g(\theovec,c)^\top H(\theovec,c)^{-1} g(\theovec,c)
  \;-\; \tfrac{1}{2}\log\det H(\theovec,c) \;+\; \mathrm{const},
  \label{eq:cioc}
\end{equation}
with $g(\theovec,c) = B(c)\theovec = \nabla_x J_\theovec(\xdemo,c)$ and
$H(\theovec,c) = \sum_k \theovec_k \nabla^2_{xx}\phi_k(\xdemo,c)$. Both terms
of \eqref{eq:cioc} are evaluated only at $\xdemo$; CIOC still performs zero
forward solves. The first term is $G$-residual-like (a Hessian-weighted
version of \eqref{eq:kkt}); the log-det term is the Gaussian normalizer and is
what prevents the fit from degenerating toward a flat cost.

\subsection{Rollout-based (bilevel) estimators}

The alternative is to actually solve the inner problem and compare the
resulting trajectory to the demonstration:
\begin{equation}
  \hat\theovec_{\mathrm{roll}}
  \;=\; \arg\min_{\theovec \in \Delta^{K-1}} \;
  \frac{1}{M}\sum_{i=1}^M \ell\big(x^\star(\theovec, c_i),\, \xdemo_i\big),
  \label{eq:bilevel}
\end{equation}
for a trajectory-space loss $\ell$ (e.g.\ squared error). The outer gradient
of \eqref{eq:bilevel} requires $\partial x^\star / \partial \theovec$, obtained
by implicit differentiation of the root condition $F(x^\star,\theovec,c)=0$:
\begin{equation}
  \frac{\partial x^\star}{\partial \theovec}
  \;=\; -\Big(\frac{\partial F}{\partial x}\Big)^{-1}
  \frac{\partial F}{\partial \theovec}
  \;=\; -H(\theovec,c)^{-1} B(c),
  \label{eq:adjoint}
\end{equation}
where the adjoint solve reuses the same $H$, $B$ that appear in
\eqref{eq:kkt}--\eqref{eq:cioc}. Equation \eqref{eq:adjoint} is only valid at
a point where $F(x,\theovec,c) = 0$ holds \emph{exactly for the $\theovec$
being evaluated} — i.e., at $x = x^\star(\theovec,c)$, not at an arbitrary
point such as $\xdemo$ for $\theovec \neq \theovec^\star$. This is the formal
reason "implicit-differentiate the demonstration directly" is not a
well-posed shortcut: for $\theovec \ne \theovec^\star$ (which includes every
$\theovec$ visited during outer optimization before convergence),
$F(\xdemo, \theovec, c) \ne 0$, so $\eqref{eq:adjoint}$ does not describe a
valid sensitivity there. Equation \eqref{eq:bilevel} avoids the issue by
always evaluating $\ell$ at a true root $x^\star(\theovec, c_i)$, paying for a
forward solve (and, since $F$ is generally nonconvex, potentially multiple
solves from different initializations) at every $\theovec$ visited.

\section{Identifiability is shared, and is a property of $\mathcal{D}$}

\begin{proposition}[Null-space ambiguity]
Let $v \in \R^K$ satisfy $B(c_i)\, v = 0$ for all $i = 1,\dots,M$, i.e.\
$v \in \bigcap_i \ker B(c_i) = \ker G$. Then for any $\theovec$ and any
$\epsilon$ small enough that $\theovec + \epsilon v$ remains feasible,
\begin{equation}
  F(\xdemo_i,\, \theovec + \epsilon v,\, c_i) = F(\xdemo_i,\theovec,c_i)
  \quad \text{for all } i,
\end{equation}
so \eqref{eq:kkt} and \eqref{eq:cioc} cannot distinguish $\theovec$ from
$\theovec + \epsilon v$: both attain the same objective value. If in addition
$x^\star(\theovec,c) = x^\star(\theovec+\epsilon v, c)$ for all $c$ in the
support of $\mathcal{D}$ (true whenever the inner optimum is a function only
of the stationarity condition along directions $v$ leaves invariant),
\eqref{eq:bilevel} is equally unable to distinguish them.
\end{proposition}

\begin{remark}
$\lambda_{\min}(G)$, normalized by $\mathrm{tr}(G)$, quantifies how close
$\mathcal{D}$ is to admitting such a $v$: $\lambda_{\min}(G) \approx 0$ means
some direction of $\Delta^{K-1}$ is (near-)unidentifiable from $\mathcal{D}$
regardless of estimator. $G$ is accumulated over the $M$ demonstration
contexts (\S2.1), so $\lambda_{\min}(G)$ is monotonically improvable only by
adding contexts that excite new feature directions — not by any property of
the estimator applied to a fixed $\mathcal{D}$. This is why $M$, not solver
sophistication, is the axis that governs recoverability, and why it governs
\eqref{eq:kkt}, \eqref{eq:cioc}, and \eqref{eq:bilevel} identically.
\end{remark}

\section{Where the two families diverge: model mismatch}

The estimators differ not in identifiability but in what happens when the
generative assumption behind \eqref{eq:kkt}/\eqref{eq:cioc} is violated,
i.e.\ when $\xdemo_i \ne x^\star(\theovec^\star, c_i)$ exactly — either
because of injected demonstration noise, or because $\xdemo_i$ is itself the
output of a forward solve carried out only to finite numerical precision.

\subsection{Analytic estimators: unmodeled residual is aliased into $\theovec$}

Write $\xdemo_i = x^\star(\theovec^\star,c_i) + \delta_i$ for a perturbation
$\delta_i$ (demonstration noise, or numerical residual from an
imperfectly-converged forward solve). Then
\begin{equation}
  B(c_i)\,\theovec^\star \Big|_{\xdemo_i}
  \;=\; \underbrace{B(c_i)\,\theovec^\star\Big|_{x^\star(\theovec^\star,c_i)}}_{=\,0}
  \;+\; \big(\nabla_x B(c_i)\,\theovec^\star\big)\,\delta_i \;+\; O(\|\delta_i\|^2)
  \;\ne\; 0,
\end{equation}
so \eqref{eq:kkt} is minimized by some $\hat\theovec \ne \theovec^\star$ that
partially cancels the $\delta_i$-induced residual against the identifiable
directions of $\theovec$, since the objective contains no term that can
attribute residual to $\delta_i$ rather than to $\theovec$. There is no
mechanism in \eqref{eq:kkt} or \eqref{eq:cioc} to separate "$\delta_i$ large"
from "$\theovec$ wrong": both estimators output a point estimate with no
internal signal that the generative assumption has failed. This is a
\emph{silent} failure mode: the returned $\hat\theovec$ carries no lower
confidence at large $\|\delta\|$ than at $\delta = 0$.

\begin{remark}[Numerical precision as unmodeled $\delta$]
If $\xdemo_i$ is produced by a forward solve converged only to tolerance
$\tau$ (e.g.\ $\tau \sim 10^{-6}$ under IEEE float32, versus
$\tau \sim 10^{-15}$ under float64), then $\|\delta_i\| = O(\tau)$ even when
no demonstration noise is injected ($\sigma = 0$ in the notation of \S5). This
$\delta_i$ is then indistinguishable, from \eqref{eq:kkt}'s point of view,
from genuine demonstration noise of the same magnitude: the nominal
$\sigma = 0$ operating point is not actually exact, but exact up to the
solver's numerical floor. Reducing $\tau$ (float64 forward solves) removes
this unmodeled contribution to $\delta_i$ and restores the theoretical
$\sigma = 0$ guarantee that \eqref{eq:kkt} is unbiased and free.
\end{remark}

\subsection{Rollout-based estimators: the loss is evaluated on a rolled-out trajectory}

Equation \eqref{eq:bilevel} evaluates $\ell$ between $\xdemo_i$ and
$x^\star(\theovec, c_i)$, a trajectory that is re-optimized (not merely
locally probed) for each candidate $\theovec$. A perturbation $\delta_i$ on
the demonstration therefore enters only through $\ell$, which — unlike the
KKT residual — is not required to vanish at the true $\theovec^\star$; it is
free to trade a nonzero loss against a $\theovec$ estimate that still fits the
bulk of the trajectory. This is what produces graceful degradation instead of
aliasing: the estimator can report "the fit is imperfect" (nonzero $\ell$)
rather than silently absorbing the imperfection into $\theovec$.

The same mechanism exposes basis misspecification. If the true generator
$\theovec^\star$ is defined over a feature basis $\phi^{\mathrm{true}}$ not
spanned by the fitting basis $\phi$ (e.g.\ dynamic-cost demonstrations fit
against a kinematic-cost $\phi$), then no $\theovec \in \Delta^{K-1}$ makes
$B(c)\theovec = 0$ at the demonstrations in general, but \eqref{eq:kkt} still
returns \emph{some} minimizer with no indication that the residual floor is
structural rather than noise-induced. Equation \eqref{eq:bilevel}, by
contrast, reports the regret $\ell(x^\star(\hat\theovec,c),\xdemo)$
explicitly, which is large and visible under basis misspecification.

\section{Summary}

\begin{itemize}
\item Identifiability ($\lambda_{\min}(G) > 0$) is a property of the
  demonstration set $\mathcal{D}$ — specifically, of how many independent
  feature directions its contexts excite — and constrains
  \eqref{eq:kkt}, \eqref{eq:cioc}, and \eqref{eq:bilevel} identically. No
  estimator recovers more than $\mathcal{D}$ identifies.
\item Given identifiability, the estimators differ only in how they handle
  $\xdemo_i \ne x^\star(\theovec^\star,c_i)$: analytic estimators
  (\S2.1) have no mechanism to detect this and alias the discrepancy into
  $\theovec$ (a confident, silent failure); rollout-based estimators (\S2.2)
  pay for a forward solve per candidate $\theovec$ but report the mismatch as
  loss/regret rather than hiding it in the estimate.
\item Numerical solver precision is a special case of the same
  $\delta_i \ne 0$ mechanism: an under-converged (e.g.\ float32) forward
  solve manufactures an unmodeled residual at nominal $\sigma = 0$ that
  analytic estimators cannot distinguish from genuine demonstration noise.
  Tightening the forward-solve tolerance (float64) removes this artifact and
  restores the $\sigma=0$ exactness that \eqref{eq:kkt}/\eqref{eq:cioc} are
  theoretically entitled to.
\end{itemize}

\end{document}
